\chapter{Характерный \texorpdfstring{$C_4$}{C4}-селектор и допустимость аффинного стока}
\label{ch:v6-spectral-transition-compacton-c4-affine-selector-admissibility-gate}

\section{Постановка после провала динамического захвата}

Предыдущий гейт показал, что при $\kappa=2\pi$ компакттон принадлежит
непрерывному семейству точных орбит $F^2=-1$, а состояния с фазами
$\pm i$ являются лишь двумя одношаговыми собственными ветвями. Однако в
Томе V уже построены проекторы характеров четырёхтактного цикла, а в
первой части Тома VI найден аффинный сток ранга три. Поэтому перед
замораживанием ветви необходимо проверить единственную оставшуюся
структурную зацепку:

\begin{quote}
может ли характер $C_4$ выделить пару $\pm i$ и передать это различие
аффинному стоку без выбора четырёхцикла, оси или диссипативного закона
вручную?
\end{quote}

Антикруговая сверка использует
\path{version5_transition_primitive_scientific_language_gate},
\path{version4_affine_family_carrier_gate},
\path{version6_existing_multiplicity_resonant_sink_gate},
\path{version6_single_thread_global_cycle_sewing_gate} и
\path{version6_spectral_transition_discrete_compacton_dynamical_capture_gate}.

\section{Точный кинематический селектор существует}

На точном двухузловом компакттонном многообразии один шаг сводится на
паре направленных блоков к оператору
\begin{equation}
 U_4=
 \begin{pmatrix}
  0&-1\\
  1&0
 \end{pmatrix},
 \qquad
 U_4^2=-I,
 \qquad
 U_4^4=I.
 \label{eq:v6-compacton-reduced-c4-step}
\end{equation}
Внутренний сбалансированный вектор переносится тождественно, поэтому
полный редуцированный шаг равен $U_4\otimes I_3$.

Канонические проекторы двух присутствующих характеров имеют вид
\begin{equation}
 P_{\pm i}=\frac14\sum_{n=0}^3(\pm i)^{-n}U_4^n,
 \qquad
 \rank P_{\pm i}=1.
 \label{eq:v6-compacton-pm-i-character-projectors}
\end{equation}
Для нормированного состояния $\Psi$ обозначим
\begin{equation}
 w_{\pm}=\langle\Psi,(P_{\pm i}\otimes I_3)\Psi\rangle.
\end{equation}
Тогда $w_++w_-=1$ и
\begin{equation}
 \mathcal D_\chi(\Psi)
 :=1-\left|\langle\Psi,(U_4\otimes I_3)\Psi\rangle\right|^2
 =4w_+w_-\geq0.
 \label{eq:v6-compacton-character-purity-defect}
\end{equation}

\begin{theorem}[Характерная чистота компакттона]
На точном многообразии $F^2=-1$ нуль функционала
\eqref{eq:v6-compacton-character-purity-defect} состоит ровно из двух
ветвей
\begin{equation}
 F\Psi=+i\Psi,
 \qquad
 F\Psi=-i\Psi.
\end{equation}
В блочной записи это условие эквивалентно
$v_R=\pm i v_L$. Следовательно, один коэффициент-свободный неотрицательный
инвариант одновременно фиксирует равенство двух узловых норм,
коллинеарность внутренних векторов и относительную фазу $\pm\pi/2$.
\end{theorem}

Скан $64$ равноотстоящих фаз дал нули только при
$\varphi/\pi=1/2,3/2$. На $64$ случайных комплексных состояниях тождество
$\mathcal D_\chi=4w_+w_-$ выполнено с максимальным остатком
$8.89\cdot10^{-16}$.

Это сильнее прежнего наблюдения о двух специальных фазах: теперь имеется
точный диагностический селектор. Но пока он является функцией состояния,
а не членом выведенного действия и не законом эволюции.

\section{Почему аффинный триплет выглядит подходящим}

Аффинный носитель имеет каноническое разложение
\begin{equation}
 \mathbb C^4=\operatorname{im}P_1\oplus\operatorname{im}P_3,
 \qquad
 P_1=\frac14J_4,
 \qquad
 P_3=I-P_1.
\end{equation}
Если выбрать четырёхцикл $c=(1234)$, его спектр на перестановочном
носителе равен
\begin{equation}
 \Spec(c)=\{1,i,-1,-i\}.
\end{equation}
Равномерный синглет несёт характер $1$, а внутри $P_3\mathbb C^4$
действительно находятся три одномерных сектора $i,-1,-i$. Поэтому
аффинный сток содержит ровно два фазовых канала, совпадающих с двумя
компакттонными ветвями.

Однако полная симметрия родителя равна
$\operatorname{AGL}(2,2)\simeq S_4$. В ней имеется шесть ориентированных
четырёхциклов, образующих три различных подгруппы $C_4$. Полная симметрия
не выбирает одну из них. Обращение ориентации дополнительно меняет местами
характеры $i$ и $-i$.

\begin{proposition}[Неканоничность аффинного $C_4$]
Выбранный четырёхцикл разлагает аффинный триплет на три характера, но это
разложение не является $S_4$-инвариантным. Три неориентированные плоскости
$P_i+P_{-i}$ и шесть ориентированных проекторов $P_i$ образуют нетривиальные
орбиты полного $S_4$. Поэтому назначение компакттонного $C_4$ аффинным
меткам вводит новое дискретное нарушение симметрии.
\end{proposition}

Машинный аудит воспроизвёл коммутант полного перестановочного
представления:
\begin{equation}
 \operatorname{End}_{S_4}(\mathbb C^4)=\operatorname{span}\{I,J\},
\end{equation}
а после ограничения на $P_3$ его комплексная размерность равна единице.
Это именно лемма Шура: любой сохраняющий полный аффинный родитель оператор
на триплете скалярен.

\section{Каноническая стрелка не очищает характер}

Связь между аффинным и семейным триплетами равна
\begin{equation}
 X=\rho V,
 \qquad
 VV^*=I_3,
 \qquad
 V^*V=P_3.
\end{equation}
После отождествления двух триплетов она пропорциональна тождеству. Поэтому
она не различает $i,-1,-i$ до выбора конкретного четырёхцикла.

Более того, любой линейный унитарный шаг, коммутирующий с $U_4$, сохраняет
веса $w_+$ и $w_-$. Из
\eqref{eq:v6-compacton-character-purity-defect} немедленно следует
\begin{equation}
 \mathcal D_\chi(F\Psi)=\mathcal D_\chi(\Psi).
\end{equation}
Следовательно, $C_4$-эквивариантная когерентная динамика распознаёт
характеры, но не очищает их смесь до одного из двух секторов.

Конечный прямосуммовой сток также не создаёт необратимость. Для
минимального резонансного гамильтониана
\begin{equation}
 H_{\rm sink}=
 \begin{pmatrix}
 0&I_3\\ I_3&0
 \end{pmatrix}
\end{equation}
вероятность исходного угла равна $\cos^2t$: при $t=\pi/2$ она полностью
переходит в сток, но при $t=\pi$ точно возвращается. Это частный явный
случай квантовой рекуррентности \cite{compacton-c4-bocchieri1957}.

Необратимое приближение к нулю $\mathcal D_\chi$ требует открытой
полугруппы, измерения, шума или нелинейной обратной связи. Общая форма
марковской квантовой полугруппы известна
\cite{compacton-c4-gks1976,compacton-c4-lindblad1976}, но проект пока не
вывел её операторы скачка, резервуар и скорость из родительской
суперсвязности.

\begin{nogocriterion}[Диагностический селектор не равен динамике]
Нельзя объявлять пару $\pm i$ динамически выбранной только потому, что
для неё построен положительный функционал с точным нулевым множеством.
Необходимо вывести эволюцию, монотонно уменьшающую
$\mathcal D_\chi$, не выбирая один из трёх аффинных $C_4$-подъёмов и не
вставляя операторы Линдблада вручную.
\end{nogocriterion}

\section{Вердикт и следующий гейт}

Ретроспективный поиск дал настоящую, но кинематическую зацепку:
$\mathcal D_\chi$ без коэффициентов сводит непрерывное компакттонное
семейство к Real-сопряжённой паре $\pm i$. Аффинный триплет после выбора
четырёхцикла содержит соответствующие принимающие каналы.

Одновременно два независимых запрета закрывают заявленный механизм:
полный $S_4$ не выбирает подгруппу $C_4$, а конечная когерентная связь не
создаёт сжатия и сохраняет характерные веса. Поэтому найден новый точный
наблюдаемый дефект, но не эндогенный trigger, скорость или захват.

Следующий гейт
\path{version6_spectral_transition_discrete_compacton_c4_boundary_eta_dissipation_gate}
проверяет последний явно сформулированный вариант: выделенная ось,
граничное нарушение $S_4\to C_4$ и слабая диссипация из уже известной
эта/Pfaffian-фазы.

Нулевая гипотеза: даже после разрешённого граничного выбора
эта/Pfaffian-фаза ориентирует характеры, но не выводит положительную
скорость и внедиагональный оператор захвата. Стоп-критерий: если для
сжатия требуется свободный оператор Линдблада или коэффициент, после
этого теста ветвь замораживается.

\begin{protocol}
\begin{enumerate}
 \item редуцированный закон $U_4^2=-I$, $U_4^4=I$: \textbf{точный};
 \item проекторы характеров $\pm i$: \textbf{построены};
 \item коэффициент-свободный дефект $\mathcal D_\chi$: \textbf{построен};
 \item его нули: \textbf{ровно две ветви $\pm i$};
 \item выравнивание двух узлов внутри точного многообразия:
 \textbf{следует из нуля $\mathcal D_\chi$};
 \item четырёхциклы полного $S_4$: \textbf{шесть};
 \item различные подгруппы $C_4$: \textbf{три};
 \item канонический выбор одной подгруппы: \textbf{нет};
 \item коммутант на аффинном триплете: \textbf{только скаляры};
 \item каноническая связь $\rho V$ очищает характер: \textbf{нет};
 \item конечный когерентный сток является аттрактором: \textbf{нет};
 \item R2, эндогенный trigger: \textbf{не закрыт};
 \item R3, скорость: \textbf{не закрыта};
 \item R4, единственный Real-парный endpoint: \textbf{кинематически выделен,
 динамически не выбран}.
\end{enumerate}
\end{protocol}

Машинный сертификат сохранён в
\path{s2t_v6_spectral_transition_compacton_c4_affine_selector_admissibility_gate_results.json}.

\begin{thebibliography}{9}
\bibitem{compacton-c4-bocchieri1957}
P.~Bocchieri and A.~Loinger,
\emph{Quantum Recurrence Theorem},
Physical Review 107 (1957), 337--338,
\path{doi:10.1103/PhysRev.107.337}.

\bibitem{compacton-c4-gks1976}
V.~Gorini, A.~Kossakowski, and E.~C.~G.~Sudarshan,
\emph{Completely positive dynamical semigroups of N-level systems},
Journal of Mathematical Physics 17 (1976), 821--825,
\path{doi:10.1063/1.522979}.

\bibitem{compacton-c4-lindblad1976}
G.~Lindblad,
\emph{On the generators of quantum dynamical semigroups},
Communications in Mathematical Physics 48 (1976), 119--130,
\path{doi:10.1007/BF01608499}.
\end{thebibliography}